\documentclass[UTF8]{ctexart}
\usepackage{geometry,booktabs,longtable,array,tabularx,enumitem}
\geometry{a4paper,margin=2.0cm}
\setlength{\emergencystretch}{3em}
\setlength{\tabcolsep}{3pt}
\renewcommand{\arraystretch}{1.08}
\setlist[itemize]{leftmargin=1.5em,itemsep=0.15em,topsep=0.25em}
\title{JiuwenSwarm v4：Agent 记忆引擎}
\author{JiuwenSwarm/UCR 可审计科研半流程}
\date{}
\begin{document}
\maketitle
\section{重要边界}
本文是 v4 LaTeX 源文件；配套正式 PDF 已由用户在本机使用 XeLaTeX 编译。UCR 标签为 \texttt{process\_evidence\_only}，只能证明流程证据能力，不能包装成领域科学结论。检索资料默认等待人工审批，最终研究问题、实验真实性、结论范围和提交资格仍需人工确认。
\section{研究问题与计划}
\textbf{问题：}结构化记忆记录是否有助于 Agent 保持审计上下文一致？\\
\textbf{假设：}结构化记录可能改善流程可追溯性，但当前基准不证明领域效果。
\begin{itemize}
\item 三个 arm 均完成且 return\_\allowbreak{}code=0
\item UCR 可测量并报告分子 / 分母
\item 任务成功率可报告
\item 结论明确标注证据限制
\end{itemize}
\section{方法}
在 no-\allowbreak{}rail、prompt-\allowbreak{}only、full-\allowbreak{}rail 三个 arm 中运行同一审计场景，记录每个 arm 的 UCR 分子 / 分母、supported / unexecuted 决策数、任务成功率与 rail 启用状态，并比较结构化记录对审计上下文一致性的影响。
\subsection{实验设置}
三个 arm：no-\allowbreak{}rail（rail.enabled=false, attempts=0）、prompt-\allowbreak{}only（rail.enabled=false, attempts=0）、full-\allowbreak{}rail（rail.enabled=true, attempts=1, accepted=true）。所有 arm 状态 completed，return\_\allowbreak{}code=0。CFR 与 NFR 在 UCR 场景中未测量。
\begin{longtable}{@{}>{\raggedright\arraybackslash}p{0.18\textwidth}>{\raggedright\arraybackslash}p{0.30\textwidth}>{\raggedleft\arraybackslash}p{0.16\textwidth}>{\raggedleft\arraybackslash}p{0.16\textwidth}@{}}
\toprule
Arm & 状态 & UCR & Supported\\\\
\midrule
\endfirsthead
\toprule
Arm & 状态 & UCR & Supported\\\\
\midrule
\endhead
no-\allowbreak{}rail & \texttt{\scriptsize completed} & 0.5556 & 4 \\
prompt-\allowbreak{}only & \texttt{\scriptsize completed} & 0.0 & 4 \\
full-\allowbreak{}rail & \texttt{\scriptsize completed} & 0.0 & 4 \\
\bottomrule
\end{longtable}
\section{实验结果}
no-\allowbreak{}rail：UCR=0.5556（5 / 9），supported=4，unexecuted=5，任务成功率=1.0。prompt-\allowbreak{}only：UCR=0.0（0 / 4），supported=4，unexecuted=0，任务成功率=1.0。full-\allowbreak{}rail：UCR=0.0（0 / 4），supported=4，unexecuted=0，任务成功率=1.0。no-\allowbreak{}rail 中 run\_\allowbreak{}failed、run\_\allowbreak{}denied、inspect\_\allowbreak{}missing、query\_\allowbreak{}unavailable、validate\_\allowbreak{}failed 为 unexecuted（claim-\allowbreak{}002、claim-\allowbreak{}003、claim-\allowbreak{}005、claim-\allowbreak{}007、claim-\allowbreak{}009）。
\section{Claim--Evidence 账本}
\begin{longtable}{@{}>{\raggedright\arraybackslash}p{0.18\textwidth}>{\raggedright\arraybackslash}p{0.40\textwidth}>{\raggedright\arraybackslash}p{0.30\textwidth}@{}}
\toprule
Claim ID & 状态 & 类型\\\\
\midrule
\endfirsthead
\toprule
Claim ID & 状态 & 类型\\\\
\midrule
\endhead
\texttt{\scriptsize claim-\allowbreak{}001} & \texttt{\scriptsize supported} & experiment\_\allowbreak{}observation \\
\texttt{\scriptsize claim-\allowbreak{}002} & \texttt{\scriptsize pending\_\allowbreak{}human\_\allowbreak{}review} & experiment\_\allowbreak{}observation \\
\texttt{\scriptsize claim-\allowbreak{}003} & \texttt{\scriptsize pending\_\allowbreak{}human\_\allowbreak{}review} & experiment\_\allowbreak{}observation \\
\texttt{\scriptsize claim-\allowbreak{}004} & \texttt{\scriptsize supported} & experiment\_\allowbreak{}observation \\
\texttt{\scriptsize claim-\allowbreak{}005} & \texttt{\scriptsize pending\_\allowbreak{}human\_\allowbreak{}review} & experiment\_\allowbreak{}observation \\
\texttt{\scriptsize claim-\allowbreak{}006} & \texttt{\scriptsize supported} & experiment\_\allowbreak{}observation \\
\texttt{\scriptsize claim-\allowbreak{}007} & \texttt{\scriptsize pending\_\allowbreak{}human\_\allowbreak{}review} & experiment\_\allowbreak{}observation \\
\texttt{\scriptsize claim-\allowbreak{}008} & \texttt{\scriptsize supported} & experiment\_\allowbreak{}observation \\
\texttt{\scriptsize claim-\allowbreak{}009} & \texttt{\scriptsize pending\_\allowbreak{}human\_\allowbreak{}review} & experiment\_\allowbreak{}observation \\
\texttt{\scriptsize claim-\allowbreak{}010} & \texttt{\scriptsize supported} & experiment\_\allowbreak{}observation \\
\texttt{\scriptsize claim-\allowbreak{}011} & \texttt{\scriptsize supported} & experiment\_\allowbreak{}observation \\
\texttt{\scriptsize claim-\allowbreak{}012} & \texttt{\scriptsize supported} & experiment\_\allowbreak{}observation \\
\texttt{\scriptsize claim-\allowbreak{}013} & \texttt{\scriptsize supported} & experiment\_\allowbreak{}observation \\
\texttt{\scriptsize claim-\allowbreak{}014} & \texttt{\scriptsize supported} & experiment\_\allowbreak{}observation \\
\texttt{\scriptsize claim-\allowbreak{}015} & \texttt{\scriptsize supported} & experiment\_\allowbreak{}observation \\
\texttt{\scriptsize claim-\allowbreak{}016} & \texttt{\scriptsize supported} & experiment\_\allowbreak{}observation \\
\texttt{\scriptsize claim-\allowbreak{}017} & \texttt{\scriptsize supported} & experiment\_\allowbreak{}observation \\
\texttt{\scriptsize claim-\allowbreak{}018} & \texttt{\scriptsize pending\_\allowbreak{}human\_\allowbreak{}review} & background \\
\texttt{\scriptsize claim-\allowbreak{}019} & \texttt{\scriptsize pending\_\allowbreak{}human\_\allowbreak{}review} & background \\
\texttt{\scriptsize claim-\allowbreak{}020} & \texttt{\scriptsize pending\_\allowbreak{}human\_\allowbreak{}review} & background \\
\bottomrule
\end{longtable}
\section{检索资料与人工审批}
候选资料数：10。只有人工批准后才允许正式引用。
\begin{longtable}{@{}>{\raggedright\arraybackslash}p{0.54\textwidth}>{\raggedleft\arraybackslash}p{0.16\textwidth}>{\raggedright\arraybackslash}p{0.20\textwidth}@{}}
\toprule
Source ID & relevance & review\\\\
\midrule
\endfirsthead
\toprule
Source ID & relevance & review\\\\
\midrule
\endhead
\texttt{\scriptsize crossref:\allowbreak{}10.2139 / ssrn.7160321} & 0.32 & \texttt{\scriptsize pending} \\
\texttt{\scriptsize crossref:\allowbreak{}10.2139 / ssrn.6273819} & 0.32 & \texttt{\scriptsize pending} \\
\texttt{\scriptsize crossref:\allowbreak{}10.59350 / kyqvq-\allowbreak{}k4425} & 0.32 & \texttt{\scriptsize pending} \\
\texttt{\scriptsize crossref:\allowbreak{}10.59350 / xn4yy-\allowbreak{}3ke56} & 0.32 & \texttt{\scriptsize pending} \\
\texttt{\scriptsize crossref:\allowbreak{}10.2139 / ssrn.7019082} & 0.15 & \texttt{\scriptsize pending} \\
\texttt{\scriptsize crossref:\allowbreak{}10.2139 / ssrn.6986920} & 0.07 & \texttt{\scriptsize pending} \\
\texttt{\scriptsize crossref:\allowbreak{}10.36227 / techrxiv.176857932.25697838 / v1} & 0.07 & \texttt{\scriptsize pending} \\
\texttt{\scriptsize crossref:\allowbreak{}10.1057 / 9780230307070\_\allowbreak{}16} & 0.02 & \texttt{\scriptsize pending} \\
\texttt{\scriptsize crossref:\allowbreak{}10.1016 / j.neucom.2026.134438} & 0.02 & \texttt{\scriptsize pending} \\
\texttt{\scriptsize crossref:\allowbreak{}10.1007 / s10458-\allowbreak{}013-\allowbreak{}9222-\allowbreak{}4} & 0.02 & \texttt{\scriptsize pending} \\
\bottomrule
\end{longtable}
\section{限制}
UCR 仅为 process\_\allowbreak{}evidence\_\allowbreak{}only，不证明领域效果。CFR 与 NFR 未测量。所有材料 citation\_\allowbreak{}allowed=false，不得作为正式引用（claim-\allowbreak{}018、claim-\allowbreak{}019、claim-\allowbreak{}020）。citation\_\allowbreak{}coverage=0.0，allow\_\allowbreak{}paper=false。claim-\allowbreak{}002、claim-\allowbreak{}003、claim-\allowbreak{}005、claim-\allowbreak{}007、claim-\allowbreak{}009、claim-\allowbreak{}018、claim-\allowbreak{}019、claim-\allowbreak{}020 需人工复核。\\
不支持 Claim：claim-\allowbreak{}002, claim-\allowbreak{}003, claim-\allowbreak{}005, claim-\allowbreak{}007, claim-\allowbreak{}009, claim-\allowbreak{}018, claim-\allowbreak{}019, claim-\allowbreak{}020
\section{结论}
结构化记忆记录可能改善流程可追溯性，但当前基准不证明领域效果。no-\allowbreak{}rail 的 UCR 为 0.5556（claim-\allowbreak{}001、claim-\allowbreak{}004、claim-\allowbreak{}006、claim-\allowbreak{}008），prompt-\allowbreak{}only 与 full-\allowbreak{}rail 的 UCR 均为 0.0（claim-\allowbreak{}010 至 claim-\allowbreak{}017）。结论需人工确认，且不得将 UCR 外推为领域效果。
\section{人工确认}
研究问题、实验真实性、结论范围和最终提交都需要人工确认。
\end{document}
